% !PRINTED 289
\weberpage{289}
%<<CONT
\[
\begin{array}{ll}
1. & \tfrac{1}{2}\big[E(f; \varphi, \psi) - A(f; \varphi, \psi)\big], \\
2. & \tfrac{1}{2}\big[E(\varphi; \psi, f) - A(\varphi; \psi, f)\big], \\
3. & \tfrac{1}{2}\big[E(\psi; f, \varphi) - A(\psi; f, \varphi)\big], \\
4. & \tfrac{1}{2}\big[E(f; \psi, \varphi) - A(f; \psi, \varphi)\big], \\
5. & \tfrac{1}{2}\big[E(\psi; \varphi, f) - A(\psi; \varphi, f)\big], \\
6. & \tfrac{1}{2}\big[E(\varphi; f, \psi) - A(\varphi; f, \psi)\big],
\end{array}
\]
\gesperrt{und es gilt nun der fundamentale Satz, dass diese sechs Bestimmungen
dieselbe Zahl ergeben.}

Um ihn zu beweisen, durchlaufen wir die Curve $f$ in positivem
Sinne, und achten auf die sämmtlichen Schnittpunkte mit $\varphi$
und $\psi$. Der Weg längs der Curve $f$ wird dabei ebenso oft in
den Binnenraum $(\varphi, \psi)$ eintreten müssen, wie er aus ihm heraustritt,
weil er wieder in seinen Ausgangspunkt zurückkehren
muss. Die Curve $f$ tritt aber ein in den Punkten $E(f; \varphi, \psi)$
und $A(f; \psi, \varphi)$ und tritt aus in den Punkten $A(f; \varphi, \psi)$ und
$E(f; \psi, \varphi)$. Also ist
\[
E(f; \varphi, \psi) + A(f; \psi, \varphi) = A(f; \varphi, \psi) + E(f; \psi, \varphi),
\]
wodurch die Uebereinstimmung von 1. und 4. nachgewiesen ist.

\weberfigureimage{weber-b1-fig09.png}{0.42}

Wenn wir nun zweitens die Curve $\varphi$ in negativem Sinne
durchlaufen, und auf ihre Schnittpunkte mit $f$ achten, so ergiebt
sich, dass jeder Punkt $E(f; \varphi, \psi)$ zugleich ein Punkt $E(\varphi; f, \psi)$,
und jeder Punkt $A(f; \varphi, \psi)$ ein Punkt $A(\varphi; f, \psi)$ ist. Denn in
der Fig.~9, worin der Punkt $\alpha$ irgend einen der Schnittpunkte
von $f$ und $\varphi$ darstellt, ist dieser Punkt ein $E(f; \varphi, \psi)$ und
ein $E(\varphi; f, \psi)$, wenn $\psi$ in $\alpha$ positiv ist, und ein $A(f; \varphi, \psi)$ und
ein $A(\varphi; f, \psi)$, wenn $\psi$ in $\alpha$ negativ ist.

Daraus folgt die Uebereinstimmung von 1. mit 6. Es folgt
ferner durch nochmalige Anwendung des ersten Schlusses die
Uebereinstimmung von 6. mit 2. und dann durch Anwendung des
zweiten Schlusses die von 2. mit 5. und endlich des ersten die
von 5. mit 3., \gesperrt{also die Uebereinstimmung aller sechs Ausdrücke.}

Man ist daher berechtigt, die so bestimmte Zahl schlechtweg
als die \gesperrt{Charakteristik des Functionensystems} $(f, \varphi, \psi)$ zu
bezeichnen.
%CONT>>
